% GPU 显存占用分解图 —— 堆叠柱状图
% 展示不同模型/精度组合下：权重(静态) + KV Cache(@8K) + 框架开销
\documentclass[border=5pt,tikz]{standalone}
\usepackage{pgfplots}
\pgfplotsset{compat=1.18}
\usepackage{xeCJK}
\setCJKmainfont{STSong}

\begin{document}
\begin{tikzpicture}
\begin{axis}[
    ybar stacked,
    width=18cm, height=9cm,
    ylabel={显存占用 (GB)},
    xlabel={模型配置},
    xtick=data,
    xticklabels={
        7B FP16, 7B INT4,
        14B FP16, 14B INT4,
        32B FP16, 32B INT8, 32B INT4,
        70B FP16, 70B INT4
    },
    xticklabel style={rotate=35, anchor=east, font=\small},
    ymin=0, ymax=150,
    bar width=0.7cm,
    legend style={at={(0.02,0.98)}, anchor=north west, font=\small},
    grid=major,
    title={GPU 显存占用分解（8K 上下文）},
    title style={font=\bfseries\large},
    % Draw GPU VRAM limit lines
    extra y ticks={16,24,32,48,80},
    extra y tick labels={16GB,24GB,32GB,48GB,80GB},
    extra y tick style={grid=major, grid style={dashed, red!50, line width=0.5pt}},
]

% 模型权重 (静态)
\addplot[fill=blue!60] coordinates {
    (0,14.0)(1,3.5)(2,28.0)(3,7.0)(4,64.0)(5,32.0)(6,16.0)(7,140.0)(8,35.0)
};
\addlegendentry{模型权重}

% KV Cache @ 8K
\addplot[fill=orange!70] coordinates {
    (0,4.3)(1,4.3)(2,6.5)(3,6.5)(4,10.0)(5,10.0)(6,10.0)(7,18.0)(8,18.0)
};
\addlegendentry{KV Cache (8K)}

% 推理框架开销
\addplot[fill=gray!40] coordinates {
    (0,2.0)(1,2.0)(2,2.0)(3,2.0)(4,2.5)(5,2.5)(6,2.5)(7,3.0)(8,3.0)
};
\addlegendentry{框架开销}

% GPU 显存上限标注
\draw[red!70, thick, dashed] (axis cs:-0.5,24) -- (axis cs:8.5,24)
    node[pos=0.08, above, font=\tiny\bfseries, red!70] {RTX 4090 24GB};
\draw[red!70, thick, dashed] (axis cs:-0.5,32) -- (axis cs:8.5,32)
    node[pos=0.08, above, font=\tiny\bfseries, red!70] {RTX 5090 32GB};
\draw[red!70, thick, dashed] (axis cs:-0.5,48) -- (axis cs:8.5,48)
    node[pos=0.08, above, font=\tiny\bfseries, red!70] {A6000 48GB};
\draw[red!70, thick, dashed] (axis cs:-0.5,80) -- (axis cs:8.5,80)
    node[pos=0.08, above, font=\tiny\bfseries, red!70] {A100 80GB};

\end{axis}

% 注解：超出 24GB 的红色标记
\node[anchor=north west, font=\footnotesize, red!80!black] at (4.2,-1.8)
    {\textbf{红色虚线} = 常见 GPU 显存上限，柱体超过虚线则需降级精度或换卡};

\end{tikzpicture}
\end{document}
